% Source: https://artofproblemsolving.com/community/c3h1248370_circle_problem
% Stable identity: aops-topic-1248370-q01
Let the chord have length $L$ and distance $d$ from the centre. The perpendicular bisector passes through the centre, so the segment from the chord to the far side of the circle has length $d+2$. Also $\left(\dfrac {L}{2}\right)^2+d^2=4$.

\begin{asy}
usepackage("amsmath");
size(3.5cm);
defaultpen(fontsize(10pt));
draw(scale(2)*unitcircle); pair Pm=(0,-1.2); real h=1.6;
draw((-h,-1.2)--(h,-1.2),blue); draw((0,-1.2)--(0,2),red); dot((0,0)); label("$O$",(0,0),E);
label("$L$",(0,0.4),W,red); label("$L$",(0.8,-1.2),S,blue); draw((0,-1.2)--(-h,-1.2),blue+linewidth(1));
draw((0,0)--(h,-1.2),dashed); label("$2$",(h/2,-0.6),NE);
label("\mbox{}",(min(currentpicture,user=true)+max(currentpicture,user=true))/2,black); // keeps the text after the figure black when LaTeX typesets the labels inline
\end{asy}

With $d=L-2$:
\[\dfrac{L^2}{4}+(L-2)^2=4\quad\Longrightarrow\quad\dfrac{5}{4}L^2-4L=0\quad\Longrightarrow\quad L=\boxed{\dfrac{16}{5}}\]
(then $d=\dfrac{6}{5}$, and indeed $\left(\dfrac{8}{5}\right)^2+\left(\dfrac{6}{5}\right)^2=4$).
